\documentclass[11pt,a4paper]{article}
\usepackage[T1]{fontenc}
\usepackage[utf8]{inputenc}
\usepackage{lmodern}
\usepackage[margin=27mm,headheight=14pt]{geometry}
\usepackage{amsmath,amssymb,amsthm,mathtools}
\usepackage{microtype,booktabs,array,longtable}
\usepackage{enumitem,xurl,fancyhdr}
\usepackage[hidelinks]{hyperref}
\usepackage[nameinlink,noabbrev]{cleveref}
\hypersetup{pdftitle={Anchored Recovery of the Dyadic Uniform Spectrum},pdfauthor={Research manuscript prepared with ChatGPT},pdfsubject={Sharp stretched-exponential inverse moduli, geometric separation, and prefix stability}}
\pagestyle{fancy}\fancyhf{}
\fancyhead[L]{\small Anchored dyadic recovery}
\fancyhead[R]{\small Proofs and reproducible checks}
\fancyfoot[C]{\thepage}
\setlength{\parskip}{3pt}
\setlength{\emergencystretch}{3em}
\setlist[enumerate]{itemsep=5pt,topsep=5pt}
\numberwithin{equation}{section}
\newtheorem{theorem}{Theorem}[section]
\newtheorem{lemma}[theorem]{Lemma}
\newtheorem{proposition}[theorem]{Proposition}
\newtheorem{corollary}[theorem]{Corollary}
\theoremstyle{definition}
\newtheorem{definition}[theorem]{Definition}
\newtheorem{question}{Research question}
\theoremstyle{remark}
\newtheorem{remark}[theorem]{Remark}
% ed. (2026-09-29): unnumbered environment for editorial notes.
\newtheorem*{ednote}{Editorial note (ProveIt, 2026-09-29)}
% ed. (2026-09-30): a second unnumbered environment for the notes of 2026-09-30.
\newtheorem*{ednotelater}{Editorial note (ProveIt, 2026-09-30)}
\DeclareMathOperator{\TV}{TV}
\DeclareMathOperator{\Law}{Law}
\DeclareMathOperator{\Var}{Var}
\DeclareMathOperator{\sinc}{sinc}
\DeclareMathOperator{\supp}{supp}
\newcommand{\R}{\mathbb R}
\newcommand{\N}{\mathbb N}
\newcommand{\E}{\mathbb E}
\newcommand{\PP}{\mathbb P}
\newcommand{\A}{\mathcal A}
\newcommand{\K}{\mathcal K}
\newcommand{\eps}{\varepsilon}
\newcommand{\dd}{\,\mathrm d}
\newcommand{\norm}[1]{\left\lVert #1\right\rVert}
\newcommand{\up}{\operatorname{up}}
\newcommand{\repoSHA}{39472967ce67566d7c530fd5b8d6a3a95e68fe22}
\newcommand{\distK}{d_{\mathrm K}}

\title{\textbf{Anchored Recovery of the\\Dyadic Uniform Spectrum}\\[7pt]
\large A sharp stretched-exponential modulus,\\lacunary counterexamples to H\"older stability,\\and Lipschitz recovery of every fixed prefix}
\author{Research manuscript prepared for Vladimir Reshetnikov\\
\small AI-assisted mathematical draft with conventional proofs}
\date{29 September 2026}

\begin{document}
\maketitle
\begin{abstract}
Let $\mu_a$ be the law of $\sum_{j\ge1}a_jU_j$, where the $U_j$ are
independent uniforms on $[-1,1]$ and $a$ is a decreasing summable sequence.
We resolve the anchored dyadic recovery question left in the ProveIt manuscript
\emph{Recovering Uniform Factors}: one of the two laws being compared is now
exactly the Rvachev law $\mu_0$, with spectrum $a_j^0=2^{-j}$.
For either total variation or Kolmogorov distance, the anchored inverse modulus
has the scale
\[
 \omega(\eps)=\exp\!\left[-\Theta\!\left(\sqrt{\log(1/\eps)}\right)\right].
\]
Thus no positive H\"older exponent is possible, but recovery at the reference
law is faster than every negative power of $\log(1/\eps)$.
The lower bound holds in the predecessor's exact fixed-variance,
uniformly smooth class, and even after imposing
$a_{j+1}\le\rho a_j$ for any fixed $\rho>1/2$.
In contrast, every fixed leading prefix is locally Lipschitz-stable, with the
rest of the spectrum completely unknown subject to a support bound.
The upper proof reconstructs primitive factor periods from stable zero counts
by divisor inversion. The lower proof perturbs the constant coefficient of a
polynomial with geometrically spaced positive roots, matches many cumulants,
and separates low-frequency cancellation from high-frequency smoothing.
Explicit finite-error formulas, a testing sample-complexity theorem,
exact rational regression tests, a proof audit, and further questions are included.
The exponent $1/2$ is sharp; the leading constant in the stretched exponential
is not determined.
\end{abstract}

\noindent\textbf{Status.} This is a proposed research contribution supported by
proofs written below, not a Lean- or Rocq-checked result. The repository
question and the comparison class are identified precisely. Literature priority
has not been established exhaustively, and independent review remains necessary.

\clearpage
\begingroup
\small
\setlength{\parskip}{0pt}
\tableofcontents
\endgroup
\clearpage

\section{The question, the comparison class, and the contribution}
\label{sec:question}

\subsection{An anchored question that pairwise instability does not answer}

The ProveIt Fabius-function research collection contains the manuscript
\emph{Recovering Uniform Factors: Exact Moment Fibres and Logarithmic
Instability Near the Fabius--Rvachev Law} \cite{predecessor}.
Its second research question asks whether a power-law inverse estimate might
hold relative to the \emph{single exact dyadic reference spectrum}.
The fourth asks whether geometric separation improves recovery, and the fifth
asks about recovering only a fixed number of leading factors.
These are logically distinct from its pairwise construction: two laws can be
extremely close to each other and both approach a reference law without either
being comparably close to that reference law.

Here every lower-bound comparison is directly against the exact law $\mu_0$.
Moreover, the perturbations do not insert a cluster of almost equal scales.
They change finitely many existing even-indexed dyadic scales by small relative
amounts and retain the odd-indexed scales exactly. Consequently they remain
uniformly geometrically separated. This is the main structural difference
from the predecessor's construction.

The inspected repository reference is
\begin{quote}\small\ttfamily\raggedright
\repoSHA.
\end{quote}
The exact source path and the relevant question locations are recorded in
\cref{sec:provenance}. The repository itself labels the predecessor as an
unreviewed archival arrival. We do not use its quantitative theorems as
unproved inputs: the present construction and both analytic bounds are proved
independently below. Its class definition is retained deliberately, so that the
new obstruction cannot be dismissed as a change to rougher densities.

\subsection{Definitions}

For $L>0$ let
\begin{equation}
 \A_L=\left\{a_1\ge a_2\ge\cdots\ge0:\ \sum_{j\ge1}a_j\le L\right\},
 \qquad
 X_a=\sum_{j\ge1}a_jU_j,\qquad \mu_a=\Law(X_a).
 \label{eq:class}
\end{equation}
The series converges absolutely for every outcome. Zero entries are padding;
they are not extra identifiable factors. Define
\begin{equation}
 d_\infty(a,b)=\sup_{j\ge1}|a_j-b_j|,
 \qquad a^0=(2^{-1},2^{-2},\ldots),\qquad \mu_0=\mu_{a^0}.
\end{equation}
We write $f_0=\up$ for the density of $\mu_0$. With this convention,
$\supp\mu_0=[-1,1]$, $\Var(X_{a^0})=1/9$, and $f_0(0)=1$.

For probability measures let
\begin{align}
 \TV(\mu,\nu)&=\sup_B|\mu(B)-\nu(B)|,
 \\
 \distK(\mu,\nu)&=\sup_{x\in\R}|F_\mu(x)-F_\nu(x)|.
\end{align}
In particular $\distK\le\TV$, and for densities total variation is one-half
the $L^1$ distance. Every logarithm is natural unless a base is displayed.

\begin{definition}[The predecessor's smooth class]
\label{def:K}
Let $\K\subset\A_3$ consist of spectra whose laws have variance $1/9$ and
a density $f_a\in C_c^\infty(\R)$ satisfying
\begin{equation}
 \norm{f_a^{(r)}}_\infty
 \le 2^{r+1}\norm{f_0^{(r)}}_\infty\qquad(r=0,1,2,\ldots).
 \label{eq:K}
\end{equation}
For $1/2<\rho<1$ put
\[
 \K_\rho=\{a\in\K:a_{j+1}\le\rho a_j\text{ for every }j\}.
\]
Both classes contain $a^0$.
\end{definition}

\begin{definition}[Anchored inverse modulus]
For a class $\mathcal C$ containing $a^0$ and a law distance $D$, set
\begin{equation}
 \omega_{\mathcal C,D}(\eps)
 =\sup\{d_\infty(a,a^0):a\in\mathcal C,
                       \ D(\mu_a,\mu_0)\le\eps\}.
 \label{eq:modulus}
\end{equation}
This is an anchored modulus, not a supremum over two freely varying spectra.
\end{definition}

\subsection{Main results}

\begin{theorem}[Sharp anchored scale]
\label{thm:main}
Take $D=\TV$ or $D=\distK$. For each of the classes
\[
 \mathcal C=\A_L\ (L>1),\qquad \mathcal C=\K,
 \qquad\text{or}\qquad \mathcal C=\K_\rho\ (1/2<\rho<1),
\]
there are positive constants $c,C,\eps_0$ such that
\begin{equation}
 \exp\!\left[-C\sqrt{\log(1/\eps)}\right]
 \le\omega_{\mathcal C,D}(\eps)
 \le\exp\!\left[-c\sqrt{\log(1/\eps)}\right]
 \quad(0<\eps<\eps_0).
 \label{eq:main}
\end{equation}
More quantitatively, the proofs give
\begin{equation}
 \sqrt{\log2}
 \le \liminf_{\eps\downarrow0}
 \frac{\log(1/\omega_{\mathcal C,D}(\eps))}{\sqrt{\log(1/\eps)}}
 \le \limsup_{\eps\downarrow0}
 \frac{\log(1/\omega_{\mathcal C,D}(\eps))}{\sqrt{\log(1/\eps)}}
 \le \frac{20}{3}\sqrt{\log2}.
 \label{eq:constant-bracket}
\end{equation}
No assertion that the two limiting constants coincide is made.
\end{theorem}

\begin{theorem}[Fixed prefixes are Lipschitz-stable]
\label{thm:prefix-intro}
For every $L\ge1$ and every fixed integer $r\ge1$ there are
$C_{L,r}<\infty$ and $\eps_{L,r}>0$ such that
\begin{equation}
 \max_{1\le j\le r}|a_j-2^{-j}|
 \le C_{L,r}\distK(\mu_a,\mu_0)
 \quad\text{if }a\in\A_L,\quad
 \distK(\mu_a,\mu_0)\le\eps_{L,r}.
 \label{eq:prefix-main}
\end{equation}
The same is true with total variation. No tail separation, fixed tail, or
finite-capacity assumption is required. One can take
$\log C_{L,r}\le(\log2)r^2+O_L(r+1)$.
\end{theorem}

There is no contradiction. A fixed prefix and the whole infinite spectrum are
different losses. The constants in the prefix theorem deteriorate rapidly as
the prefix grows; the least detectable perturbed coordinate in the lower-bound
construction also moves to infinity.

The statistical consequence is likewise stated in a form that avoids a
pointwise-estimation trap. Testing the fixed null $a=a^0$ against
$d_\infty(a,a^0)\ge\delta$ has sample complexity
\begin{equation}
 N_*(\delta)=\exp\!\left[\Theta\!\left(\log^2(1/\delta)\right)\right]
 \label{eq:testing-intro}
\end{equation}
for fixed nontrivial error levels. We prove this for alternatives in $\K_\rho$
and give an upper test valid on all of $\A_3$.

\subsection{What is new here, and what is established background}

Qualitative identifiability of generalized uniform sums is established work:
Billey and Swanson give a unique decreasing-sequence parametrization and a
topological theory that also permits a Gaussian component \cite{billey}.
Arias de Reyna is a primary source for the compactly supported smooth function
and its transform conventions \cite{arias}. The use of Newton identities,
Rouch\'e's theorem, Plancherel's theorem, and empirical distribution-function
bounds is classical.

The proposed contribution is the specific quantitative combination: an anchored
stretched-exponential modulus with matching exponent; a counterexample family
inside the exact smooth class and within any separation class $\rho>1/2$;
and Lipschitz recovery of a fixed prefix despite an unrestricted infinite tail.
These are not presented as new versions of identifiability itself. Nor does
this article settle the predecessor's \emph{pairwise} optimal logarithmic
modulus, its exact-support-endpoint question, or general minimax estimation over
all spectra.

\section{Analytic and algebraic preliminaries}
\label{sec:prelim}

\subsection{Entire transforms and their zeros}

\begin{lemma}\label{lem:transform}
For $a\in\A_L$ the characteristic function extends to the entire function
\begin{equation}
 \phi_a(z)=\E e^{izX_a}=\prod_{j\ge1}\sinc(a_jz),
 \qquad \sinc z=\frac{\sin z}{z},\quad\sinc0=1.
 \label{eq:product}
\end{equation}
Its positive zeros, with multiplicity, are exactly
\begin{equation}
 \left\{\frac{k\pi}{a_j}:a_j>0,\ k\ge1\right\}.
 \label{eq:zero-multiset}
\end{equation}
In particular it has no nonreal zeros and no extra zeros created by the
infinite product.
\end{lemma}
\begin{proof}
Since $|X_a|\le L$, the expectation is entire by locally uniform domination.
Independence gives the finite products, and absolute convergence of the random
series gives their locally uniform convergence to the expectation. Alternatively,
$\sum_j|\sinc(a_jz)-1|$ converges uniformly on every compact set because
$\sum a_j^2<\infty$. After removing finitely many initial factors, every
remaining factor is close to one and its analytic logarithm is summable.
Thus the tail product is nonzero at every point not already a factor zero.
Only finitely many factors can vanish at a given nonzero point, because such a
factor must have $a_j\ge\pi/|z|$. This proves the multiplicity assertion.
\end{proof}

For the dyadic law put
\begin{equation}
 \Phi_0(w)=\phi_0(2\pi w)
 =\prod_{j\ge1}\sinc(\pi 2^{1-j}w).
 \label{eq:normalized-transform}
\end{equation}
Its positive zero at the integer $m$ has multiplicity
\begin{equation}
 c_m=1+v_2(m).
 \label{eq:multiplicity}
\end{equation}
Indeed the contributing factors are precisely those with $2^{j-1}\mid m$.

\subsection{Two measurement norms control the same contours}

\begin{lemma}\label{lem:measurement}
Suppose $\mu,\nu$ are supported in $[-L,L]$. Then
\begin{align}
 |\phi_\mu(z)-\phi_\nu(z)|
 &\le2e^{L|\operatorname{Im}z|}\TV(\mu,\nu),
 \label{eq:tv-entire}\\
 |\phi_\mu(z)-\phi_\nu(z)|
 &\le2L|z|e^{L|\operatorname{Im}z|}\distK(\mu,\nu).
 \label{eq:ks-entire}
\end{align}
\end{lemma}
\begin{proof}
The first inequality follows by integrating $e^{izx}$ against the signed
measure $\mu-\nu$, whose total mass variation is $2\TV(\mu,\nu)$.
For the second, Stieltjes integration by parts gives
\[
 \phi_\mu(z)-\phi_\nu(z)
 =-iz\int_{-L}^{L}e^{izx}(F_\mu(x)-F_\nu(x))\dd x.
\]
One may first integrate over a slightly larger interval, where the boundary
terms vanish, so atoms at the endpoints cause no difficulty. Taking absolute
values gives the claim.
\end{proof}

\subsection{Power sums and the logarithm of sinc}

For $|u|<\pi$, Euler's sine product and absolute summation give
\begin{equation}
 \log\sinc u
 =-\sum_{k\ge1}\frac{\zeta(2k)}{k\pi^{2k}}u^{2k}.
 \label{eq:log-sinc}
\end{equation}
For real $|u|<\pi$ this is the real logarithm of a positive number. In the
complex disk it is the analytic branch vanishing at zero. Thus equality of
$\sum a_j^{2k}$ for $1\le k\le m$ implies equality of cumulants, and hence
moments, through degree $2m+1$, since all odd cumulants vanish. We shall use
this only for finite blocks, where all algebraic manipulations are immediate.

\section{Recovering primitive periods from perturbed zero counts}
\label{sec:periods}

The crucial discrete step is not generic stability of a multiple root.
It uses the fact that each uniform factor contributes \emph{an entire harmonic
progression} of zeros.

\begin{lemma}[Stable divisor reconstruction]
\label{lem:periods}
Fix $\eta=1/10$ and an integer $M\ge4$. Let
$0<\beta_1\le\beta_2\le\cdots$ be a finite or locally finite sequence,
with $\beta_i\to\infty$ in the infinite case. Suppose the multiset
$\{k\beta_i:k\ge1\}$ has the following properties up to $2M+\eta$:
all its points lie in the intervals
\[
 I_m=(m-\eta,m+\eta),\qquad 1\le m\le2M,
\]
and exactly $1+v_2(m)$ points, counting multiplicity, lie in $I_m$.
Then for every power of two $d\le M$ there is exactly one primitive period
$\beta(d)\in I_d$, and there are no other primitive periods at most
$M+\eta$. Moreover,
\begin{equation}
 |\beta(d)-d|\le\frac{\eta\beta(d)}{M}.
 \label{eq:period-accuracy}
\end{equation}
If $b_i=1/(2\beta_i)$, decreasingly ordered and padded by zeros when necessary,
then
\begin{equation}
 d_\infty(b,a^0)\le\frac1{2M}.
 \label{eq:period-spectral}
\end{equation}
\end{lemma}

\begin{proof}
Any primitive period $\beta\le M+\eta$ itself is an observed zero, so it lies
in a unique $I_d$, $1\le d\le M$. Write $\beta=d+e$, $|e|<\eta$.
For all $k$ with $k\beta\le2M+\eta$, let $m_k$ be the integer satisfying
$|k\beta-m_k|<\eta$. Taking $m_0=0$, consecutive integers satisfy
\[
 |(m_k-m_{k-1})-d|<3\eta<1.
\]
Therefore $m_k-m_{k-1}=d$, and induction gives
\begin{equation}
 m_k=kd,\qquad |ke|<\eta.
 \label{eq:no-phase-slip}
\end{equation}
This prevents a progression from gradually drifting to a different lattice
class. Merely locating its first zero would not give this conclusion.

If $kd\le M$, then $k\beta\le M+M\eta<2M+\eta$, so this zero is indeed
observed and belongs to $I_{kd}$. Conversely, a zero in $I_m$ for $m\le M$
comes from a primitive period at most $M+\eta$, and \eqref{eq:no-phase-slip}
shows that its associated integer $d$ divides $m$.
Writing $n_d$ for the number of primitive periods in $I_d$, we obtain
\begin{equation}
 1+v_2(m)=\sum_{d\mid m}n_d\qquad(1\le m\le M).
 \label{eq:divisor-system}
\end{equation}
This triangular system has a unique solution. The solution is $n_d=1$ when
$d$ is a power of two and $n_d=0$ otherwise, because a positive integer $m$
has exactly $1+v_2(m)$ power-of-two divisors.

For a recovered period put $K=\lfloor(2M+\eta)/\beta\rfloor$. Since
$\beta\le M+\eta$, we have $K\ge M/\beta$. Equation
\eqref{eq:no-phase-slip} yields $|e|\le\eta/K\le\eta\beta/M$.
Consequently
\[
 \left|\frac1{2\beta(d)}-\frac1{2d}\right|
 \le\frac{\eta}{2Md}.
\]
These are exactly the leading dyadic factors with $d\le M$. Every remaining
period is greater than $M+\eta$, while the next dyadic primitive period is
greater than $M$. The two nonnegative tails are therefore bounded above by
$1/(2M)$. This proves \eqref{eq:period-spectral}.
\end{proof}

\begin{remark}
The lemma is an approximate-lattice statement with exact multiplicities.
It is robust to multiplicities in the unknown spectrum and makes no a priori
separation assumption on that spectrum. Divisor inversion forces the relevant
separation after the zero counts have been established.
\end{remark}

\section{An explicit upper modulus and the prefix theorem}
\label{sec:upper}

\subsection{A uniform lower bound away from the zero lattice}

Keep $\eta=1/10$ and define, for $M\ge4$,
\begin{align}
 r_0&=2\pi\eta,& R_M&=2\pi(2M+\eta),
 &J_M&=\left\lceil\log_2(2R_M)\right\rceil,\\
 B_M&=e^{-1}\left(\frac{4\eta}{R_M}\right)^{J_M},
 &\eps_M(L)&=\frac{B_M}{4(1+LR_M)e^{Lr_0}}.
 \label{eq:thresholds}
\end{align}

\begin{lemma}\label{lem:lower-contour}
The bound $|\phi_0(z)|\ge B_M$ holds on every circle
$|z-2\pi m|=r_0$, $1\le m\le2M$, and at every real
$z\in[r_0,R_M]$ outside the corresponding open disks.
For real $0\le z\le r_0$, one has $\phi_0(z)>1/2$.
\end{lemma}
\begin{proof}
For every complex $u$,
\[
 |\sin u|^2=\sin^2(\operatorname{Re}u)+\sinh^2(\operatorname{Im}u)
 \ge\frac4{\pi^2}\operatorname{dist}(u,\pi\mathbb Z)^2.
\]
Here one reduces the real part to $[-\pi/2,\pi/2]$, uses
$|\sin x|\ge2|x|/\pi$, and uses $|\sinh y|\ge|y|$.
At any point specified in the lemma,
$\operatorname{dist}(z,2\pi\mathbb Z)\ge r_0$.
Since $2^j\pi\mathbb Z\subset2\pi\mathbb Z$ for $j\ge1$,
\[
 |\sinc(z/2^j)|
 \ge\frac{2\operatorname{dist}(z,2^j\pi\mathbb Z)}{\pi|z|}
 \ge\frac{4\eta}{R_M}\qquad(1\le j\le J_M).
\]
For $j>J_M$, $|z|2^{-j}\le1/4$. The power series of sinc gives
$|\sinc u|\ge1-|u|^2/5\ge e^{-|u|^2}$ for $|u|\le1/2$.
Their product is at least
\[
 \exp\left[-R_M^2\sum_{j>J_M}4^{-j}\right]\ge e^{-1}.
\]
This proves the claimed product bound. Finally, for real $z$,
$\phi_0(z)=\E\cos(zX_{a^0})\ge1-z^2/18$, which exceeds $1/2$ on
$[0,r_0]$.
\end{proof}

\begin{theorem}[Finite-error recovery certificate]
\label{thm:finite-upper}
Let $L\ge1$, $M\ge4$, and $a\in\A_L$. If either
\[
 \TV(\mu_a,\mu_0)\le\eps_M(L)
 \quad\text{or}\quad
 \distK(\mu_a,\mu_0)\le\eps_M(L),
\]
then
\begin{equation}
 d_\infty(a,a^0)\le\frac1{2M}.
 \label{eq:finite-upper}
\end{equation}
\end{theorem}
\begin{proof}
By \cref{lem:measurement}, the transform difference on the circles and real
segments of \cref{lem:lower-contour} is at most $B_M/2$.
Rouch\'e's theorem therefore gives exactly $1+v_2(m)$ zeros inside each
circle centered at $2\pi m$. All zeros of the unknown transform are real
by \cref{lem:transform}. On the intervening real segments the unknown
transform cannot vanish. On $[0,r_0]$ the same conclusion follows from the
last assertion of \cref{lem:lower-contour} and the measurement bound.

In the normalized coordinate $w=z/(2\pi)$, the primitive periods of the
unknown factors are $\beta_j=1/(2a_j)$ for $a_j>0$. They satisfy all
hypotheses of \cref{lem:periods}. That lemma gives the desired bound.
\end{proof}

\begin{corollary}[Upper stretched-exponential rate]
\label{cor:upper}
For either law distance, and any fixed $L\ge1$,
\begin{equation}
 \omega_{\A_L,D}(\eps)
 \le C_L\exp\left[-\sqrt{\log2}\sqrt{\log(1/\eps)}\right]
 \quad\text{for sufficiently small }\eps.
 \label{eq:upper-asymptotic}
\end{equation}
\end{corollary}
\begin{proof}
The explicit threshold satisfies
\begin{equation}
 \log\frac1{\eps_M(L)}
 =\frac{(\log M)^2}{\log2}+O_L(\log M+1).
 \label{eq:threshold-asymptotic}
\end{equation}
To check this, use $R_M\asymp M$ and
$J_M=\log M/\log2+O(1)$ in \eqref{eq:thresholds}.
If $s=\log(1/\eps)$, take
$M=\lfloor\exp(\sqrt{s\log2}-A_L)\rfloor$, with $A_L$ sufficiently
large. The upper bound implicit in \eqref{eq:threshold-asymptotic} then
ensures $\eps\le\eps_M(L)$. Apply \cref{thm:finite-upper}; rounding $M$
changes only the multiplicative constant.
\end{proof}

\subsection{Why multiple transform zeros do not destroy prefix Lipschitzness}

\begin{proof}[Proof of \cref{thm:prefix-intro}]
Fix $r$ and take $M=\max(4,2^{r-1})$. When the law distance
$\delta\le\eps_M(L)$, the proof of \cref{thm:finite-upper} identifies
exactly one primitive period $\beta_j$ close to $d_j=2^{j-1}$ for each
$j\le r$. Write $\Phi_a(w)=\phi_a(2\pi w)$.

Let $S_j(a)$ be the sum, with multiplicities, of the zeros in
$|w-d_j|<\eta$. The progression identification in \cref{lem:periods} gives
\begin{equation}
 S_j(a)=\sum_{i=1}^{j}2^{j-i}\beta_i,
 \qquad \beta_j=S_j(a)-2S_{j-1}(a),\qquad S_0(a)=0.
 \label{eq:centroid-recursion}
\end{equation}
The key is to estimate the \emph{sum} of a multiple-root cluster before
attempting to identify its individual roots.

On the contour $C_j:|w-d_j|=\eta$ we have
$|\Phi_a/\Phi_0-1|\le1/2$. The principal logarithm of this ratio is
well defined on a neighborhood of the contour. The argument principle and
integration by parts along the closed contour yield
\begin{align}
 S_j(a)-S_j(a^0)
 &=\frac1{2\pi i}\int_{C_j}w
       \left(\frac{\Phi_a'}{\Phi_a}-\frac{\Phi_0'}{\Phi_0}\right)\dd w
 \\
 &=-\frac1{2\pi i}\int_{C_j}\log(\Phi_a/\Phi_0)\dd w.
 \label{eq:centroid-log}
\end{align}
No logarithm inside the full disk is asserted: a branch on an annular
neighborhood of the contour suffices for this identity.
Using $|\log(1+u)|\le2|u|$ for $|u|\le1/2$, we get
\begin{equation}
 |S_j(a)-S_j(a^0)|
 \le\frac{4\eta(1+LR_M)e^{Lr_0}}{B_M}\,\delta.
 \label{eq:centroid-bound}
\end{equation}
The recursion \eqref{eq:centroid-recursion} bounds $|\beta_j-d_j|$ by
three times this quantity. Since $\beta_j\ge(1-\eta)d_j$,
\[
 |a_j-2^{-j}|=\frac{|\beta_j-d_j|}{2\beta_jd_j}
 \le\frac{2(1+LR_M)e^{Lr_0}}{B_M}\,\delta.
\]
Thus one may take $\eps_{L,r}=\eps_M(L)$ and the displayed coefficient
as $C_{L,r}$. Its logarithm is $(\log2)r^2+O_L(r+1)$ by the same
calculation as \eqref{eq:threshold-asymptotic}.
\end{proof}

\begin{remark}[A distinction from generic polynomial root bounds]
A generic $j$-fold zero can split on a scale $\delta^{1/j}$. Here its
constituent zeros are not arbitrary: they come from already identified
harmonic progressions. Cluster sums vary linearly with the perturbation,
and the recursion removes the earlier progressions. This extra structure
is exactly what restores exponent one for every fixed prefix.
\end{remark}

\section{A positive-root perturbation of existing dyadic factors}
\label{sec:construction}

\subsection{A constant-coefficient perturbation with controlled real roots}

Fix $0<\tau\le10^{-3}$, and let $n\ge2$. Define
\begin{equation}
 \begin{aligned}
 x_j&=16^{-j}\quad(1\le j\le n),&
 P_n(x)&=\prod_{j=1}^n(x-x_j),\\
 \Delta_n&=\tau\prod_{j=1}^nx_j=\tau16^{-n(n+1)/2},&
 Q_n&=P_n-\Delta_n.
 \end{aligned}
 \label{eq:root-polynomials}
\end{equation}
Only the constant coefficient changes. This elementary operation will preserve
many power sums without requiring two moving reference spectra.

\begin{lemma}[Separated positive roots]\label{lem:root-perturbation}
The polynomial $Q_n$ has exactly one simple root $y_j$ in each interval
$(3x_j/4,5x_j/4)$. In particular $y_1>\cdots>y_n>0$.
For $\ell=n-j$, put
\begin{equation}
 w_\ell=(4/3)^\ell16^{-\ell(\ell+1)/2}.
 \label{eq:w}
\end{equation}
Then
\begin{align}
 \frac{|y_j-x_j|}{x_j}&\le2\tau w_{n-j},
 &\left|\frac{\sqrt{y_j}}{\sqrt{x_j}}-1\right|
       &\le2\tau w_{n-j},                                      \label{eq:root-relative}\\
 |\sqrt{y_n}-\sqrt{x_n}|&\ge\frac\tau3\,4^{-n}. \label{eq:root-lower}
\end{align}
The power sums satisfy
\begin{equation}
 \sum_{j=1}^ny_j^k=\sum_{j=1}^nx_j^k\quad(1\le k<n),
 \qquad
 \sum_{j=1}^ny_j^n-\sum_{j=1}^nx_j^n=n\Delta_n.
 \label{eq:newton-match}
\end{equation}
\end{lemma}

\begin{proof}
Write $q=1/16$. On the entire interval $[3x_j/4,5x_j/4]$, the product
of distances to the preceding roots is bounded below by
\[
 \prod_{i<j}|x-x_i|
 \ge\prod_{i<j}x_i\prod_{k\ge1}(1-\tfrac54q^k)
 \ge\frac9{10}\prod_{i<j}x_i.
\]
The last inequality follows from
$\prod(1-t_k)\ge1-\sum t_k$ for $0\le t_k\le1$ and
$(5/4)\sum_{k\ge1}q^k=1/12$.
The distances to the following roots satisfy, with $\ell=n-j$,
\[
 \prod_{i>j}|x-x_i|
 \ge(\tfrac34 x_j)^\ell
       \prod_{k\ge1}(1-\tfrac43q^k)
 \ge\frac9{10}(\tfrac34)^\ell x_j^\ell,
\]
since $(4/3)\sum_{k\ge1}q^k=4/45$.
At either endpoint of the interval, the remaining factor has magnitude
$x_j/4$, so
\begin{equation}
 |P_n(x)|
 \ge\frac{81}{400}(\tfrac34)^\ell
             16^{\ell(\ell+1)/2}\prod_{i=1}^nx_i
 \ge\frac{81}{400}\prod_{i=1}^nx_i>\Delta_n.
 \label{eq:bracket-sign}
\end{equation}
The signs at the two endpoints are opposite. Subtracting $\Delta_n$ preserves
these signs. The $n$ intervals are disjoint, hence contain $n$ different real
roots of a polynomial of degree $n$. This also proves simplicity and exhausts
all its roots.

At $y_j$, the equation $|P_n(y_j)|=\Delta_n$ and the same estimates with
the $j$th factor omitted give
\[
 \frac{|y_j-x_j|}{x_j}
 \le\frac{100}{81}\tau(\tfrac43)^\ell16^{-\ell(\ell+1)/2}
 \le2\tau w_\ell.
\]
For positive $u$, $|\sqrt{u}-1|\le|u-1|$, proving the second relative bound.
For the last root, $0<y_n<x_i$ when $i<n$, and therefore
\[
 \Delta_n=|y_n-x_n|\prod_{i<n}(x_i-y_n)
 \le |y_n-x_n|\prod_{i<n}x_i.
\]
It follows that $|y_n-x_n|\ge\tau x_n$. Dividing by
$\sqrt{y_n}+\sqrt{x_n}<3\sqrt{x_n}$ proves \eqref{eq:root-lower}.
Finally, Newton's identities show that the first $n-1$ power sums depend
only on the unchanged coefficients. In degree $n$ the constant term enters
with coefficient $n$, giving exactly the final formula in
\eqref{eq:newton-match}.
\end{proof}

\subsection{The full infinite spectrum}

Define $b^{(n)}$ by replacing just the first $n$ even-indexed factors:
\begin{equation}
 b^{(n)}_{2j}=\sqrt{y_j}\quad(1\le j\le n),
 \qquad b^{(n)}_i=2^{-i}\quad\text{for every other }i.
 \label{eq:b-spectrum}
\end{equation}
The order of the factors does not change. Indeed, the odd factor immediately
before the $j$th replacement is $2\cdot4^{-j}$, the replacement lies between
$\sqrt{3/4}\,4^{-j}$ and $\sqrt{5/4}\,4^{-j}$, and the next odd factor is
$4^{-j}/2$.

\begin{proposition}[Size, variance, and moment matching]
\label{prop:construction-size}
The spectra \eqref{eq:b-spectrum} satisfy
\begin{equation}
 \frac\tau3\,4^{-n}
 \le d_\infty(b^{(n)},a^0)
 \le\norm{b^{(n)}-a^0}_1
 \le3\tau4^{-n}.
 \label{eq:b-gap}
\end{equation}
Their laws have exactly the same moments as $\mu_0$ through degree $2n-1$,
and have a different moment in degree $2n$. In particular their variance is
exactly $1/9$. All their laws are supported in $[-2,2]$, and
$b^{(n)}\in\A_L$ eventually for every $L>1$.
\end{proposition}
\begin{proof}
The lower bound is \eqref{eq:root-lower}. For the upper bound, observe that
$4^\ell w_\ell\le(1/3)^\ell$. Thus
\[
 \sum_{j=1}^n|\sqrt{y_j}-4^{-j}|
 \le2\tau4^{-n}\sum_{\ell=0}^{n-1}4^\ell w_\ell
 \le3\tau4^{-n}.
\]
The support assertions follow by summing the coefficients.
Equation \eqref{eq:newton-match}, applied to their squared lengths, preserves
the full even power sums through exponent $2n-2$. The logarithmic transform
identity \eqref{eq:log-sinc} then preserves cumulants and moments through
$2n-1$. Its coefficient in degree $2n$ is nonzero, and the power-sum change
there is $n\Delta_n\ne0$, so that moment is different. The variance assertion
is the case $k=1$ of \eqref{eq:newton-match}.
\end{proof}

\subsection{Membership in the exact smooth class}

Preserving the variance alone would not answer the question in the
predecessor's class. We next check every one of its infinitely many derivative
bounds, with its stated constants.

\begin{lemma}[Derivative norms of the reference and a convolution bound]
\label{lem:derivatives}
The density $f_0$ is smooth and
\begin{equation}
 \norm{f_0^{(r)}}_\infty=2^{r(r+1)/2}\qquad(r\ge0).
 \label{eq:exact-derivative-norm}
\end{equation}
For any summable spectrum all of whose entries are positive, its density is
smooth, and
\begin{equation}
 \norm{f_a^{(r)}}_\infty
 \le\frac1{2\prod_{i=1}^{r+1}a_i}\qquad(r\ge0).
 \label{eq:general-derivative-bound}
\end{equation}
\end{lemma}
\begin{proof}
Arbitrarily many positive uniform factors imply arbitrarily fast polynomial
decay of the characteristic function, and hence a $C^\infty$ density by
Fourier inversion. The density has compact support when the spectrum is
summable. The first factor of $a^0$ has constant density one on
$[-1/2,1/2]$, and the remaining sum is supported in the same interval.
Consequently $\norm{f_0}_\infty=f_0(0)=1$.

The tail after this first factor has the law of $X_{a^0}/2$.
Differentiating its convolution with that first uniform gives
\[
 f_0'(x)=2\bigl(f_0(2x+1)-f_0(2x-1)\bigr).
\]
After $r$ further derivatives,
\[
 f_0^{(r+1)}(x)=2^{r+1}
       \bigl(f_0^{(r)}(2x+1)-f_0^{(r)}(2x-1)\bigr).
\]
The two nonzero support interiors in this formula are disjoint; each argument
ranges over the full interval $[-1,1]$ on its respective half. Smoothness
makes every boundary derivative zero. The supremum norm therefore multiplies
by exactly $2^{r+1}$ at each step, proving
\eqref{eq:exact-derivative-norm}.

For the general bound, differentiate each of the first $r$ uniform densities
once in the distributional sense. The resulting signed measures have total
variation $1/a_i$. Keep the next uniform as a density with supremum norm
$1/(2a_{r+1})$ and the rest as a probability measure. Convolution and the
already established smoothness give \eqref{eq:general-derivative-bound}.
\end{proof}

\begin{proposition}[Smoothness and geometric separation survive]
\label{prop:K-membership}
For $0<\tau\le10^{-3}$, every spectrum $b^{(n)}$ belongs to $\K$.
Given any fixed $\rho\in(1/2,1)$, the choice
\begin{equation}
 0<\tau\le\tau_\rho:=
 \min\left(10^{-3},\frac{2\rho-1}{8(2\rho+1)}\right)
 \label{eq:tau-rho}
\end{equation}
ensures $b^{(n)}\in\K_\rho$ for every $n\ge2$.
\end{proposition}
\begin{proof}
Since $w_\ell\le(1/12)^\ell$, the sum of all relative changes to the
nontrivial coordinates is at most
\[
 \sum_{j=1}^n\left|\frac{b^{(n)}_{2j}}{a^0_{2j}}-1\right|
 \le\frac{24\tau}{11}<\frac12.
\]
For every prefix, the product of its coefficient ratios is consequently at
least $1-24\tau/11>1/2$. Applying \eqref{eq:general-derivative-bound},
\[
 \norm{f_{b^{(n)}}^{(r)}}_\infty
 \le\frac1{\prod_{i=1}^{r+1}2^{-i}}
 =2^{(r+1)(r+2)/2}
 =2^{r+1}\norm{f_0^{(r)}}_\infty.
\]
The support and variance conditions were proved in
\cref{prop:construction-size}. This establishes $\K$ membership.

Every coordinate ratio $b_i^{(n)}/a_i^0$ lies in $[1-2\tau,1+2\tau]$.
Therefore
\[
 \frac{b_{i+1}^{(n)}}{b_i^{(n)}}
 \le\frac{1+2\tau}{2(1-2\tau)}.
\]
This is at most $\rho$ whenever
$\tau\le(2\rho-1)/(2(2\rho+1))$; \eqref{eq:tau-rho} is a stricter
convenient choice. The index $i=2n$, where the unchanged tail resumes,
is included in the same inequality.
\end{proof}

\section{A quantitative bound against the exact reference law}
\label{sec:fourier}

Moment matching by itself is not a bound on total variation. In particular,
a low-frequency estimate cannot simply be extended to all frequencies.
The following argument has a common smooth background, a small head
perturbation, a moment-cancelled tail, and an explicit high-frequency integral.

\subsection{The unchanged odd-indexed background}

The odd-indexed factors have characteristic function
\begin{equation}
 H(t)=\prod_{\ell\ge0}\sinc\bigl(\tfrac12\,4^{-\ell}t\bigr).
 \label{eq:odd-background}
\end{equation}
For every positive integer $K$ and $t\ne0$,
\begin{equation}
 |H(t)|\le2^{K^2}|t|^{-K}.
 \label{eq:H-bound}
\end{equation}
Indeed $|\sinc u|\le\min(1,1/|u|)$ on the real axis. Apply the second
bound to the first $K$ factors and the first bound to the others. The
reciprocal product of those $K$ half-lengths is exactly $2^{K^2}$.

There is also an unchanged even tail. More precisely, set
\[
 \Xi_n(t)=H(t)\prod_{j>n}\sinc(4^{-j}t),\quad
 A_n(t)=\prod_{j=1}^n\sinc(4^{-j}t),\quad
 B_n(t)=\prod_{j=1}^n\sinc(\sqrt{y_j}t).
\]
Then the full transform difference is
\begin{equation}
 \phi_{b^{(n)}}(t)-\phi_0(t)=\Xi_n(t)(B_n(t)-A_n(t)),
 \qquad |\Xi_n(t)|\le|H(t)|.
 \label{eq:full-difference}
\end{equation}
Keeping this common even tail in the identity is important, even though only
its modulus bound is needed.

\subsection{An explicit low-frequency estimate}

For $n\ge16$ define
\begin{align}
 h&=\lfloor n/2\rfloor,&m&=\lfloor n/8\rfloor,
 &K&=\lfloor3n/10\rfloor,&T&=2^{3n/5},                 \label{eq:cutoffs}\\
 A&=5/64,&
 D_n&=2\tau h(4/3)^n16^{-(n-h)(n-h+1)/2},              \label{eq:Dn}\\
 r_n&=\frac{\sqrt{5/4}\,4^{-(h+1)}T}{\pi}.            \label{eq:rn}
\end{align}
In particular $m<n$, $K\ge1$, and $r_n<1$. The last fact follows from
$h+1>n/2$, giving $r_n<(\sqrt{5/4}/\pi)2^{-2n/5}$.
Introduce
\begin{equation}
 E_n=TD_n+\frac{2mD_n}{A}
       \max\left(1,\frac{AT^2}{\pi^2}\right)^m
       +\frac{4n r_n^{2m+2}}{(m+1)(1-r_n^2)}.
 \label{eq:En}
\end{equation}

\begin{lemma}[Head--tail cancellation]\label{lem:low-frequency}
For real $|t|\le T$ one has
\begin{equation}
 |B_n(t)-A_n(t)|\le E_n.
 \label{eq:low-frequency}
\end{equation}
\end{lemma}
\begin{proof}
Split the two products after factor $h$. For the first $h$ factors,
\cref{lem:root-perturbation} implies both
\begin{equation}
 \sum_{j\le h}|y_j-x_j|\le D_n,
 \qquad
 \sum_{j\le h}|\sqrt{y_j}-\sqrt{x_j}|\le D_n.
 \label{eq:head-change}
\end{equation}
To see this, use $x_j,\sqrt{x_j}\le1$, $n-j\ge n-h$,
$(4/3)^{n-j}\le(4/3)^n$, and the monotonicity of
$\ell(\ell+1)/2$ for $\ell\ge0$.
On the real axis, the sinc function is Lipschitz with constant at most one:
write it as $\E e^{itU}$ and use $\E|U|\le1$. Telescoping products whose
factors have modulus at most one therefore bounds the head difference by
$TD_n$.

All arguments of the tail factors have absolute value at most $\pi r_n<\pi$.
Their products are positive, at most one, and admit the real logarithm
\eqref{eq:log-sinc}. The total power sums agree through degree $n-1$.
For $1\le k\le m$, the difference of the tail power sums is consequently
the negative of the head difference. Since $x_j,y_j\le A$ and
$|u^k-v^k|\le kA^{k-1}|u-v|$ on $[0,A]$, its magnitude is at most
$kA^{k-1}D_n$. The contribution of these $m$ terms to the logarithm
difference is at most
\[
 \sum_{k=1}^m\frac{2}{k\pi^{2k}}\,kA^{k-1}D_n T^{2k}
 \le\frac{2mD_n}{A}\max\left(1,\frac{AT^2}{\pi^2}\right)^m.
\]
Here we used $\zeta(2k)<2$.

For $k\ge m+1$ we need no cancellation. There are at most $n$ factors in
each tail, and each squared scale is bounded by
$(5/4)16^{-(h+1)}$. Thus the absolute difference of its power sums is at
most twice $n$ times that bound to the $k$th power. Summing the logarithmic
series gives
\[
 4n\sum_{k\ge m+1}\frac{r_n^{2k}}k
 \le\frac{4n r_n^{2m+2}}{(m+1)(1-r_n^2)}.
\]
Since $|e^u-e^v|\le|u-v|$ for real $u,v\le0$, the same sum bounds
the difference of the two tail products. Combining the head and tail
bounds proves \eqref{eq:low-frequency}.
\end{proof}

\subsection{The full-law error and its quadratic logarithm}

\begin{theorem}[Explicit anchored total-variation bound]
\label{thm:TV-bound}
For $n\ge16$ define
\begin{equation}
 V_n=\min\left\{1,
 \left(\frac{2TE_n^2+
       8\,2^{2K^2}T^{1-2K}/(2K-1)}{2\pi}\right)^{1/2}\right\}.
 \label{eq:Vn}
\end{equation}
Then
\begin{equation}
 \TV(\mu_{b^{(n)}},\mu_0)\le V_n,
 \qquad
 V_n\le\exp\left[-\frac{9\log2}{100}\,n^2+O_\tau(n)\right].
 \label{eq:TV-superexp}
\end{equation}
The constant implicit in $O_\tau(n)$ is independent of $n$.
\end{theorem}
\begin{proof}
On $[-T,T]$, \cref{lem:low-frequency} and $|\Xi_n|\le1$ bound the
squared Fourier integral by $2TE_n^2$. On its complement,
$|B_n-A_n|\le2$, and \eqref{eq:H-bound} gives
\[
 \int_{|t|>T}|\phi_{b^{(n)}}(t)-\phi_0(t)|^2\dd t
 \le\frac{8\,2^{2K^2}T^{1-2K}}{2K-1}.
\]
Both densities are supported in $[-2,2]$. Cauchy--Schwarz and Plancherel,
with the characteristic-function convention \eqref{eq:product}, give
\[
 \TV(\mu_{b^{(n)}},\mu_0)
 =\tfrac12\norm{f_{b^{(n)}}-f_0}_1
 \le\norm{f_{b^{(n)}}-f_0}_2
 =\left(\frac1{2\pi}\int_{\R}
       |\phi_{b^{(n)}}-\phi_0|^2\dd t\right)^{1/2}.
\]
This proves the finite formula.

For clarity, the asymptotic exponents of every term are recorded explicitly.
With $a=\log2$, the floor functions in \eqref{eq:cutoffs} give
\begin{equation}
 \begin{aligned}
 \log D_n&=-\tfrac12 a n^2+O_\tau(n),
 &\log T&=\tfrac35 a n,\\
 \log r_n&=-\tfrac25 a n+O(1),
 &m&=\tfrac18 n+O(1).
 \end{aligned}
\end{equation}
Consequently the three positive terms of $E_n$ are bounded respectively by
\begin{equation}
 \exp[-\tfrac12 a n^2+O_\tau(n)],\quad
 \exp[-\tfrac7{20}a n^2+O_\tau(n)],\quad
 \exp[-\tfrac1{10}a n^2+O(n)].
 \label{eq:E-exponents}
\end{equation}
For example, the middle term adds
$m\log(AT^2/\pi^2)=\tfrac3{20}a n^2+O(n)$ to
$\log D_n$, while the last uses
$(2m+2)\log r_n=-\tfrac1{10}a n^2+O(n)$.
The low-frequency contribution after taking the square root in
\eqref{eq:Vn} is therefore at most
$\exp[-\tfrac1{10}a n^2+O_\tau(n)]$.
For the high-frequency contribution,
\[
 2K^2\log2+(1-2K)\log T=-\tfrac9{50}a n^2+O(n).
\]
Its square root has exponent $-\tfrac9{100}a n^2+O(n)$ and is the larger
of the two bounds. Combining them proves \eqref{eq:TV-superexp}.
\end{proof}

\begin{remark}[Why both parts of the Fourier argument are necessary]
The first few changed factors are so small a perturbation that their
contribution is quadratic-exponentially small in $n$. The remaining factors
are small enough that a long truncated logarithmic expansion stays valid up
to $T=2^{3n/5}$. Beyond $T$, cancellation is no longer used: the unchanged
odd-indexed factors supply the decay. This separation avoids extrapolating
a moment expansion through the zeros of sinc.
\end{remark}

\section{Completion of the anchored theorem and its interpretation}
\label{sec:completion}

\begin{proof}[Proof of \cref{thm:main}]
The upper bound follows from \cref{cor:upper}. For $\K$ and $\K_\rho$
use their inclusion in $\A_3$. In particular it yields the first inequality
in \eqref{eq:constant-bracket}.

For the lower bound fix $\tau>0$, using \eqref{eq:tau-rho} when needed.
Put $c_0=9\log2/100$ and $s=\log(1/\eps)$. By
\cref{thm:TV-bound}, there is $C_\tau$ for which
$V_n\le\exp(-c_0n^2+C_\tau n)$ for all sufficiently large $n$.
Choose
\begin{equation}
 n=\left\lceil\frac{\sqrt{s}}{\sqrt{c_0}}+B_\tau\right\rceil,
 \label{eq:n-choice}
\end{equation}
where $B_\tau$ is a sufficiently large constant. Then $V_n\le\eps$.
The constructed spectrum belongs to $\K$ and, with the smaller choice of
$\tau$, to $\K_\rho$. It belongs to $\A_L$ for sufficiently large $n$
whenever $L>1$. Since $\distK\le\TV$, it is admissible for both law
metrics. Its parameter gap satisfies
\[
 \omega_{\mathcal C,D}(\eps)\ge\frac\tau3\,4^{-n}.
\]
Taking logarithms and using \eqref{eq:n-choice},
\[
 \log(1/\omega_{\mathcal C,D}(\eps))
 \le\frac{\log4}{\sqrt{c_0}}\sqrt{s}+O_\tau(1)
 =\frac{20}{3}\sqrt{\log2}\sqrt{s}+O_\tau(1).
\]
This proves the final inequality of \eqref{eq:constant-bracket}. The
intermediate liminf--limsup inequality is automatic. The two-sided form
\eqref{eq:main} follows after slightly weakening the constants to absorb
bounded terms.
\end{proof}

\begin{corollary}[No anchored H\"older estimate, even with lacunarity]
\label{cor:no-holder}
In each class in \cref{thm:main}, no constants $C<\infty$ and $\alpha>0$
can make
\[
 d_\infty(a,a^0)\le C D(\mu_a,\mu_0)^\alpha
\]
hold throughout a sufficiently small law neighborhood of $\mu_0$.
On the other hand, for every $p>0$,
\[
 \omega_{\mathcal C,D}(\eps)=o\bigl((\log(1/\eps))^{-p}\bigr).
\]
\end{corollary}
\begin{proof}
For the first assertion the construction gives
\[
 \frac{d_\infty(b^{(n)},a^0)}{D(\mu_{b^{(n)}},\mu_0)^\alpha}
 \ge\frac\tau3\exp\bigl(\alpha c_0n^2-O_\tau(n)\bigr)\longrightarrow\infty.
\]
The second assertion follows from the upper bound because
$e^{-c\sqrt{s}}s^p\to0$ as $s\to\infty$.
\end{proof}

\subsection{Exactly which repository questions are answered}

The anchored power-law question receives a negative answer in the very smooth
class used by the predecessor, together with a positive quantitative replacement:
the sharp exponent in the square root of the logarithm. Geometric separation
with any fixed ratio strictly greater than $1/2$ does not restore a power law
at the dyadic reference. The prefix question receives a positive answer under
a single common support bound: every fixed prefix has exponent one, even with
an unknown infinite tail.

These conclusions do not say that the full inverse has this modulus between
two arbitrary nearby spectra. The special zero lattice of the exact reference
is used on every contour in the upper proof. Nor do they determine the exact
support endpoint of the perturbed laws: their half-support lengths tend to
one, but need not equal one. Fixed variance and fixed support endpoints are
different restrictions.

\subsection{Why the two exponents are geometrically natural}

Resolving factors down to size about $1/M$ requires viewing transform zeros
at frequencies of order $M$. There are about $\log_2 M$ active dyadic
factors there, and a conservative product lower bound loses about
$\exp[-(\log M)^2/\log2]$. This yields the square root when inverted.
The construction realizes the converse scale through products of geometrically
separated root distances: their logarithms are sums of arithmetic progressions,
therefore quadratic in the number of factors.

This explanation is not substituted for a proof. The finite formulas
\eqref{eq:thresholds} and \eqref{eq:Vn} establish both directions without
an appeal to numerical fitting. They also reveal the unresolved constant gap:
the contour argument and the chosen head--tail cutoffs are not simultaneously
optimized.

\section{Testing complexity and honest confidence sets}
\label{sec:statistics}

\subsection{A testing formulation that is genuinely informative}

An estimator identically equal to $a^0$ has zero loss when the truth is exactly
$a^0$. For that reason, an assertion about ``the minimax estimation rate at a
single point'' without a neighborhood or honesty condition is misleading.
We instead consider testing a fixed null against separated alternatives.

Fix $\rho\in(1/2,1)$. Let $N_\rho(\delta)$ be the least positive integer
$N$ for which there is a possibly randomized test $0\le\psi\le1$, based
on $N$ independent observations, with
\begin{equation}
 \E_{\mu_0^{\otimes N}}\psi\le\tfrac14,
 \qquad
 \sup_{\substack{a\in\K_\rho\\d_\infty(a,a^0)\ge\delta}}
       \E_{\mu_a^{\otimes N}}(1-\psi)\le\tfrac14.
 \label{eq:test-definition}
\end{equation}
The supremum is over the alternative, not over an arbitrarily chosen estimator
at the null.

\begin{theorem}[Quasi-polynomial sample complexity in inverse accuracy]
\label{thm:testing}
For fixed $\rho\in(1/2,1)$,
\begin{equation}
 N_\rho(\delta)=
 \exp\left[\Theta_\rho\bigl(\log^2(1/\delta)\bigr)\right]
 \quad(\delta\downarrow0).
 \label{eq:testing}
\end{equation}
The constants obtained here satisfy
\begin{equation}
 \frac9{400\log2}
 \le\liminf_{\delta\downarrow0}
       \frac{\log N_\rho(\delta)}{\log^2(1/\delta)}
 \le\limsup_{\delta\downarrow0}
       \frac{\log N_\rho(\delta)}{\log^2(1/\delta)}
 \le\frac2{\log2}.
 \label{eq:testing-constants}
\end{equation}
The upper test works for all alternatives in $\A_3$, not just $\K_\rho$.
The same order holds with alternatives in $\K$ or $\A_L$ for fixed $L>1$.
\end{theorem}

\begin{proof}
Let $F_N$ be the empirical distribution function. We use the classical
Dvoretzky--Kiefer--Wolfowitz--Massart bound
\begin{equation}
 \PP_F\!\left\{\sup_x|F_N(x)-F(x)|>u\right\}
 \le2e^{-2Nu^2}\qquad(u>0).
 \label{eq:DKW}
\end{equation}
Massart proved the sharp constant two; see \cite{massart} and the account
in \cite{weidudley}. This is an external statistical input, not a new
inequality claimed here.

Take $M=\max(4,\lceil1/\delta\rceil)$ and
$\eps_*=\eps_M(3)$. By \cref{thm:finite-upper}, any alternative with
$d_\infty(a,a^0)\ge\delta$ has
$\distK(\mu_a,\mu_0)>\eps_*$, since $1/(2M)<\delta$.
Reject the null when
\begin{equation}
 \sup_x|F_N(x)-F_0(x)|>\eps_*/2.
 \label{eq:cdf-test}
\end{equation}
Under the null, \eqref{eq:DKW} bounds the rejection probability by
$2\exp(-N\eps_*^2/2)$. Under any alternative, acceptance implies an empirical
CDF error larger than $\eps_*/2$ relative to its own law, and the same bound
applies. Thus a finite upper bound is
\begin{equation}
 N_\rho(\delta)\le
 \left\lceil\frac{2\log8}{\eps_M(3)^2}\right\rceil.
 \label{eq:N-upper}
\end{equation}
Equation \eqref{eq:threshold-asymptotic} proves the upper logarithmic
constant in \eqref{eq:testing-constants}.

For the lower bound choose $\tau$ as in \eqref{eq:tau-rho} and, for small
$\delta$, put
\[
 n=\left\lfloor\frac{\log(\tau/(3\delta))}{\log4}\right\rfloor.
\]
Then $n\ge16$ eventually and $d_\infty(b^{(n)},a^0)\ge\delta$.
The elementary product-measure bound gives
\[
 \TV(\mu_{b^{(n)}}^{\otimes N},\mu_0^{\otimes N})
 \le N\TV(\mu_{b^{(n)}},\mu_0)\le NV_n.
\]
For any test, its two error probabilities at these two laws sum to at least
one minus this total variation. If $N\le1/(4V_n)$, their sum is at least
$3/4$, which is incompatible with both being at most $1/4$.
Therefore $N_\rho(\delta)>\lfloor1/(4V_n)\rfloor$.
Using \eqref{eq:TV-superexp} and
$n=\log(1/\delta)/\log4+O_\rho(1)$ yields
\[
 \log N_\rho(\delta)
 \ge\frac{9}{400\log2}\log^2(1/\delta)-O_\rho(\log(1/\delta)).
\]
The same construction is available in $\K$ and eventually in $\A_L$ for
any $L>1$; in the latter case the upper test uses $\eps_M(L)$.
\end{proof}

\subsection{Confidence sets that shrink at the reference}

Fix $L\ge1$ and a coverage error $0<\alpha<1$. Define
\begin{equation}
 t_N=\sqrt{\frac{\log(2/\alpha)}{2N}},\qquad
 \mathcal C_N=\{a\in\A_L:\sup_x|F_a(x)-F_N(x)|\le t_N\}.
 \label{eq:confidence-set}
\end{equation}
For every $a\in\A_L$, \eqref{eq:DKW} gives
$\PP_a\{a\in\mathcal C_N\}\ge1-\alpha$.
At the reference law, on an event of that probability, every candidate in
$\mathcal C_N$ satisfies $\distK(\mu_a,\mu_0)\le2t_N$.
Hence
\begin{equation}
 \sup_{a\in\mathcal C_N}d_\infty(a,a^0)
 \le\omega_{\A_L,\distK}(2t_N)
 \le C_{L,\alpha}
   \exp\left[-\sqrt{\frac{\log2}{2}}\sqrt{\log N}\right]
 \label{eq:confidence-contraction}
\end{equation}
for sufficiently large $N$. For each fixed prefix, \cref{thm:prefix-intro}
instead gives a radius at most $2C_{L,r}t_N=O_{L,r,\alpha}(N^{-1/2})$.

These are honest coverage statements over $\A_L$ with contraction at the
reference. They are not claimed as a globally optimal estimator, nor as a
finite algorithm for representing the infinite-dimensional confidence set.
Numerically implementing \eqref{eq:cdf-test} also requires validated values
of $F_0$; its statistical proof assumes that reference CDF is known exactly.
A uniform CDF approximation error can be included by enlarging the threshold,
but an implementation and its numerical-error budget are separate tasks.

\section{Reproducible checks and finite-error scales}
\label{sec:checks}

\subsection{Exact arithmetic, diagnostics, and what each establishes}

The accompanying program \texttt{code/verify.py} uses Python's rational
\texttt{Fraction} arithmetic for the finite polynomial checks. Its default
choice is $\tau=1/1000$. On the run included with this article, it completed
the following checks without an assertion failure:
\begin{center}
\begin{tabular}{@{}lr@{}}
\toprule
Exact check & Number\\
\midrule
Opposite endpoint signs for the perturbed root intervals, $2\le n\le20$ & 209\\
Unchanged power sums below degree $n$, $2\le n\le20$ & 190\\
First changed power sum equals $n\Delta_n$ & 19\\
Only the constant polynomial coefficient changes & 19\\
Divisor inversion returns the power-of-two indicator, $1\le m\le512$ & 512\\
\midrule
Total finite exact checks & 949\\
\bottomrule
\end{tabular}
\end{center}
These are regression checks, not a proof for all $n$ or a formal proof of
the analytic theorems. The universal proofs are
\cref{lem:periods,lem:root-perturbation,thm:TV-bound}.

A separate part of the program uses \texttt{mpmath} for high-precision root
bisection at $n=2,4,6,8$ and for evaluations of the explicit analytic bounds.
The root calculations used 100, 100, 158, and 242 decimal digits,
respectively. These values are floating-point diagnostics, not interval
certificates. The exact rational endpoint signs and the proved isolation
lemma are kept logically separate from the floating-point residuals.
The recorded environment was Python 3.13.5 and \texttt{mpmath} 1.3.0.

\subsection{Scale of the anchored witnesses}

\begin{table}[htbp]
\centering\small
\begin{tabular}{@{}rrr@{}}
\toprule
$n$ & Gap lower bound $(\tau/3)4^{-n}$ & TV upper bound $V_n$\\
\midrule
16  & $7.7610\times10^{-14}$  & $2.1467\times10^{-6}$\\
32  & $1.8070\times10^{-23}$  & $4.9268\times10^{-26}$\\
64  & $9.7958\times10^{-43}$  & $1.2260\times10^{-106}$\\
128 & $2.8787\times10^{-81}$  & $6.8525\times10^{-434}$\\
256 & $2.4861\times10^{-158}$ & $5.3403\times10^{-1754}$\\
512 & $1.8542\times10^{-312}$ & $9.2947\times10^{-7058}$\\
\bottomrule
\end{tabular}
\caption{Rounded evaluations of the proved formulas, with $\tau=1/1000$.
The second column is a lower bound, not a measured spectral distance.
The third is an analytic upper bound, not numerical integration of an actual
TV distance. The rounding displayed is not directed interval rounding.}
\label{tab:lower-numeric}
\end{table}

The construction is intentionally conservative. At small $n$, the Fourier
upper bound can be much larger than the spectral gap. The asymptotic
obstruction emerges from the quadratic logarithm of the law bound versus the
linear logarithm of the spectral gap. For example, the value of
$-\log V_n/n^2$ at $n=512$ is approximately $0.06199$, approaching the
proved upper-bound exponent $9\log2/100\approx0.06238$.
This concerns $V_n$ alone; it does not identify the actual TV asymptotics.

\subsection{Scale of a finite recovery certificate}

\begin{table}[htbp]
\centering\small
\begin{tabular}{@{}rrrr@{}}
\toprule
$M$ & $J_M$ & $\log_{10}\eps_M(3)$ & Guaranteed spectral error $1/(2M)$\\
\midrule
4    & 7  & $-18.7738$ & $1.2500\times10^{-1}$\\
16   & 9  & $-28.9611$ & $3.1250\times10^{-2}$\\
64   & 11 & $-41.5784$ & $7.8125\times10^{-3}$\\
256  & 13 & $-56.6109$ & $1.9531\times10^{-3}$\\
1024 & 15 & $-74.0537$ & $4.8828\times10^{-4}$\\
4096 & 17 & $-93.9054$ & $1.2207\times10^{-4}$\\
\bottomrule
\end{tabular}
\caption{Rounded evaluations of \eqref{eq:thresholds} for the support bound
$L=3$. These sufficient thresholds deliberately favor a transparent uniform
proof over good practical constants.}
\label{tab:upper-numeric}
\end{table}

The upper certificate does not recover roots by numerically searching a noisy
transform. It states an implication: a proved measurement bound below the
threshold forces the correct finite divisor pattern and hence the displayed
parameter accuracy. A practical solver would need a representation of the
observed transform, validated contour bounds, and an error budget for each
numerical operation.

\subsection{How to reproduce the package}

% ed. (2026-09-29): the program now writes to data-rerun/ by default, so a
% plain run no longer overwrites the recorded files in data/.
Run \texttt{python code/verify.py} from the package root, or invoke the
script by its absolute path from another directory. It writes
\nolinkurl{verification.json}, \nolinkurl{verification.txt},
\nolinkurl{analytic_bounds.csv}, and \nolinkurl{rouche_thresholds.csv}
to \nolinkurl{data-rerun/}, or to the directory given by
\texttt{-{}-output-dir}; the recorded copies are in \nolinkurl{data/}. The exact checks use only the Python
standard library; the diagnostic portion needs \texttt{mpmath}.
The TeX source contains the displayed tables directly, so rebuilding the
article does not require Python. Two passes of \texttt{pdflatex article.tex}
resolve its references; \texttt{latexmk -pdf article.tex} is also suitable.
No shell escape or network access is needed to compile it.

\section{Proof dependencies and repository provenance}
\label{sec:provenance}

\subsection{Claim-level dependency audit}

\begin{center}
\small
\begin{tabular}{@{}p{0.27\textwidth}p{0.65\textwidth}@{}}
\toprule
Claim & Inputs and scope\\
\midrule
Finite upper certificate & Entire sinc products, the explicit contour bound,
Rouch\'e's theorem, and the stable divisor lemma. No separation assumption on
the unknown spectrum.\\[4pt]
Prefix Lipschitz theorem & The same finite zero counts plus the contour
sum-of-zeros identity. A logarithm is used only near the contour, not across
zeros in its interior.\\[4pt]
Anchored counterexamples & Exact constant-coefficient perturbation, isolated
positive roots, Newton identities, and the unchanged odd-indexed factors.\\[4pt]
Smooth-class membership & Exact reference derivative norms and the product
bound for every prefix, as well as exact variance matching.\\[4pt]
TV bound & A real logarithm only below the first tail zero, and an independent
high-frequency product bound followed by Plancherel.\\[4pt]
Sharp exponent & The upper certificate and the explicit lower witnesses,
with integer rounding handled in both conversions.\\[4pt]
Testing result & The preceding deterministic results, product-TV contraction,
and the external DKW--Massart inequality.\\
\bottomrule
\end{tabular}
\end{center}

No mathematical conclusion is inferred solely from finite numerical tests.
No theorem from an unreviewed repository research article is imported as an
assumption. Conversely, use of standard analysis is explicit: this manuscript
does not claim a new proof of Rouch\'e, Plancherel, or the DKW inequality.

\subsection{Inspected source and exact question mapping}

The principal predecessor at the pinned commit is
\begin{quote}\small\ttfamily\raggedright
Analysis/FabiusFunction/docs/\allowbreak
semi-formalized-research-frontiers/\allowbreak
drafts/inverse-and-sampling/\allowbreak
Recovering\_Uniform\_Factors\_Fabius\_Rvachev/\allowbreak
article.tex.
\end{quote}
Its content blob identifier is
\begin{quote}\small\ttfamily\raggedright
90395cdfd50b8d8e191d2753616fc47c890a3311.
\end{quote}
The repository snapshot is \texttt{\repoSHA}. It was read through the
GitHub connector, including the opening definition of $\K$ and the final
research-question section. The following titles specify the comparison more
robustly than a line number alone:
\begin{center}
\small
\begin{tabular}{@{}>{\raggedright\arraybackslash}p{0.46\textwidth}>{\raggedright\arraybackslash}p{0.46\textwidth}@{}}
\toprule
Predecessor question & Result here\\
\midrule
Pointwise recovery of the exact dyadic spectrum &
\Cref{thm:main,cor:no-holder}: sharp stretched-exponential exponent and no
H\"older estimate in the stated smooth class.\\[4pt]
Geometrically separated classes &
\Cref{prop:K-membership,thm:main}: negative power-law answer at the dyadic
anchor for every fixed ratio $\rho>1/2$; not a full pairwise classification.\\[4pt]
Fixed leading factors versus the complete sequence &
\Cref{thm:prefix-intro}: local Lipschitz recovery of every fixed prefix with
only an $\ell^1$ support bound on the unknown full spectrum.\\[4pt]
Matching statistical upper bounds &
\Cref{thm:testing} and \eqref{eq:confidence-contraction}: matching-order
anchored testing and honest confidence contraction, not the predecessor's
global minimax problem.\\
\bottomrule
\end{tabular}
\end{center}
% ed. (2026-09-29): Lean crosswalk; the article does not cite these theorems.
\begin{ednote}
Lean crosswalk, not cited by the article. The zero multiplicity
$c_m=1+v_2(m)$ of \eqref{eq:multiplicity} is the machine-checked theorem
\path{Fabius.analyticOrderAt_rvachevFourierProduct_int} in
\path{Analysis/FabiusFunction/Lean/FabiusFunction/IntegerZeroAnalyticOrder.lean}.
That $1+v_2(m)\ge r+1$ exactly when $2^r\mid m$, the divisor count behind
\cref{lem:periods}, is
\path{Fabius.dyadicZeroMultiplicity_ge_succ_iff_pow_two_dvd} in
\path{DyadicZeroMultiplicity.lean} of the same directory. The derivative
norms \eqref{eq:exact-derivative-norm} are
\path{Fabius.isGreatest_abs_iteratedDeriv_rvachevUp}
(\path{GlobalBounds.lean}), and the variance $1/9$ is
\path{Fabius.integral_sq_mul_rvachev_eq_one_ninth}
(\path{DyadicSpecializations.lean}). Exact identifiability of dyadic
exponent sequences from zero orders, a noise-free counterpart of the
divisor inversion, is proved in \path{GeneralizedRvachevIdentifiability.lean}.
No quantitative statement of this article is formalized.
\end{ednote}

The adjacent inverse-and-sampling README was also inspected. It lists the
finite-factor, Gaussian-confounding, and flat-boundary recovery articles as
separate recent arrivals. Their fixed-capacity questions are different from
the present unknown-infinite-tail prefix result. This was a targeted source
review, not an exhaustive audit of every document in ProveIt or every paper
in the literature. No repository file was modified.

\subsection{A concrete formalization route}

A natural first module is the finite divisor-count lemma: encode the observed
zero incidences as a finite multiset, prove the no-phase-slip implication, and
solve the triangular divisor system. This module does not need probability
or Fabius analysis.

A second algebraic module can formalize the geometric root intervals,
constant-coefficient perturbation, Newton identities, and the prefix product
bound. Rational endpoint inequalities offer small finite certificate instances,
while the general estimates need elementary ordered-field and geometric-series
lemmas. Positivity and preservation of sorted rank must remain explicit.

The analytic modules should then separate three interfaces: compactly supported
law distance to contour error; entire-product zeros and argument-principle
counts; and the Fourier head--tail estimate with its $L^2$-to-TV conversion.
Finally, metric corollaries and testing statements can be assembled from these
interfaces. The sum-of-zeros proof should not be replaced by a generic root
continuity theorem, since that would lose the prefix exponent one.

This is a dependency plan, not a claim that a particular current mathlib API
already exposes every required theorem. No Lean build was run, and no uncompiled
Lean skeleton is included as apparent evidence of formal verification.

\section{Further research questions}
\label{sec:further}

\begin{question}[The exact stretched-exponential constant]
Does
\[
 \lim_{\eps\downarrow0}
 \frac{\log(1/\omega_{\mathcal C,D}(\eps))}
      {\sqrt{\log(1/\eps)}}
\]
exist, and how does it depend on the class and law metric? The present interval
$[\sqrt{\log2},(20/3)\sqrt{\log2}]$ is wide. Optimizing the retained
background, the root spacing of the replaced subsequence, and the Fourier
cutoffs is one route to a better lower construction. A sharp contour product
bound, rather than a common bound for every early factor, may improve the
upper constant independently.
\end{question}

\begin{question}[Pairwise recovery on geometrically separated classes]
Determine the optimal pairwise modulus when both unknown spectra satisfy
$a_{j+1}\le\rho a_j$. The present upper argument fixes one transform to
the exact dyadic zero lattice; it does not give a pairwise theorem uniformly
on a neighborhood. Geometric separation removes the predecessor's dense
blocks but does not remove the anchored counterexamples constructed here.
A positive theorem must therefore have a modulus no better than the anchored
stretched-exponential obstruction.
\end{question}

\begin{question}[The sharp separation boundary $\rho=1/2$]
What happens if $a_{j+1}\le a_j/2$ exactly, with variance fixed at $1/9$?
Our perturbations require slack above $1/2$. At equality, adjacent one-sided
constraints may rule out the compensating scale changes that preserve many
power sums. The case $\rho<1/2$ does not contain the dyadic reference at all,
so it is a different anchored problem rather than a strengthening of the
same one.
\end{question}
% ed. (2026-09-30): answered at the dyadic reference by a later package
\begin{ednotelater}
Answered by \path{../Lacunarity_Boundary_Geometric_Uniform_Spectra/}
(batch~68 of the repository's \path{docs/incoming/}), written against this
article and in its normalization ($U_j$ uniform on $[-1,1]$, variance
$1/9$, so $\sum_ja_j^2=1/3$). If $a_{j+1}\le a_j/2$ for every $j$, then
\[
 \sup_j|a_j-2^{-j}|^2\le\tfrac{75}{4}\bigl(\tfrac{19}{675}-\E X_a^4\bigr),
\]
and the fourth-moment deficit vanishes only at $a^0$. Hence the anchored
modulus at $\rho=1/2$ is $\Theta(\sqrt\eps)$, a power law, in total
variation and in Kolmogorov distance, with matching deleted-factor
witnesses in $\K$ whose ratios are at most $1/2$; the
stretched-exponential scale of \cref{thm:main} does not persist at the
boundary. At $\rho=1/2$ it also answers the next question negatively
(there $\sum_ja_j\le1$, with equality only at $a^0$, so no witness with
$\sum_jb_j=1$ exists), and the question on prefix constants inside its
critical cone only (local Lipschitz constants $\Theta(2^r)$). The pairwise
modulus, the exact-support question with slack, and the prefix constants
on $\A_L$ remain open.
\end{ednotelater}

\begin{question}[Exact support together with exact variance]
Can the lower witnesses be chosen with $\sum_jb_j=1$ exactly, while
retaining variance $1/9$, geometric separation, and the sharp logarithmic
exponent? The construction here controls the sum to within $3\tau4^{-n}$
but does not fix it. A two-parameter polynomial perturbation may enforce
the extra square-root sum, yet its effect on low-order cumulants and on the
quadratic Fourier exponent needs proof.
\end{question}

\begin{question}[Sharp growth of prefix conditioning constants]
Find the optimal dependence on $r$ of the local Lipschitz constant and the
radius on which it applies. We prove an upper bound
$\log C_{L,r}\le(\log2)r^2+O_L(r)$, but not a matching lower coefficient.
This is the missing quantitative transition between root-$N$ confidence
for a fixed prefix and stretched-logarithmic contraction for the complete
spectrum. Allowing $r=r(N)$ to grow is an especially concrete formulation.
\end{question}

\begin{question}[Other arithmetic geometric references]
For a reference $a_j=cq^j$, when does the full inverse have the same
square-root logarithmic exponent, and when are fixed prefixes Lipschitz?
If $q$ is the reciprocal of an integer, the zero multiplicities retain a
divisibility structure, suggesting an analogue of \cref{lem:periods}.
For a general $q$, the harmonics can be incommensurate and nearly collide.
A proof must quantify those near coincidences rather than assume the dyadic
counting argument survives unchanged.
\end{question}

\begin{question}[Square-summable spectra and Gaussian confounding]
Extend quantitative anchored recovery to the parameter space allowing
square-summable but nonsummable lengths and an unknown Gaussian component,
as in the qualitative theory of Billey--Swanson \cite{billey}. Compact
support is used decisively to turn Kolmogorov error into entire-transform
error on contours. Which tail or exponential-moment hypothesis can replace
it, and which parameter metric sees escaped Gaussian variance correctly?
\end{question}
% ed. (2026-09-29): the predecessor's corresponding question is answered, for
% the Gaussian variance, by a package filed after this article's snapshot.
\begin{ednote}
The predecessor's corresponding question, ``Gaussian components and
square-summable spectra'', is answered for the Gaussian-variance functional
by the repository article
\path{inverse-and-sampling/Gaussian_Dust_Christoffel_Recovery_Uniform_Factors/}
(filed 2026-09-29, after the snapshot this article read; unreviewed), in a
model with a known smooth compactly supported background: without a tail
restriction the exact minimax risk for the variance is $V/2$ at every
sample size. Quantitative recovery of the scales themselves in that space,
anchored or pairwise, is not treated there.
\end{ednote}

\begin{question}[Validated inversion and proof-carrying certificates]
Build a finite-prefix solver that outputs both candidate lengths and
certificates for the contour inequalities and divisor counts. A useful first
milestone is a checker for a fixed $M$ with rational enclosure data and a
proved tail estimate for the reference sinc product. The statistical CDF
error and the numerical transform error must be tracked separately. The
explicit thresholds in \eqref{eq:thresholds} are a correctness baseline,
not a recommendation to use their conservative constants unchanged.
\end{question}

\begin{question}[Adaptive confidence away from the exact anchor]
Characterize reference spectra for which honest confidence sets have
root-$N$ contraction for every fixed prefix, and quantify the whole-spectrum
contraction across different references. The confidence construction
\eqref{eq:confidence-set} is honest globally but its proved contraction rate
here is anchored. Uniform adaptation would require a comparison theorem
between local zero geometry and the inverse modulus, not merely the
pointwise analysis of $\mu_0$.
\end{question}

\section{Conclusion}

Exact identifiability of a uniform-convolution spectrum conceals several
distinct quantitative regimes. At the dyadic Rvachev law, the entire infinite
spectrum has an anchored modulus with the sharp scale
$\exp[-\Theta(\sqrt{\log(1/\eps)})]$. It admits no positive H\"older
exponent even under fixed variance, the predecessor's uniform bounds on all
derivatives, and geometric separation arbitrarily close to the dyadic ratio.
Yet every fixed leading prefix is locally Lipschitz when the unknown tail
is constrained only by a common support bound.

Both statements arise from the special structure of uniform factors. Harmonic
zero progressions allow stable divisor reconstruction and stable cluster
centroids; a constant-coefficient perturbation of separated squared scales
hides a moving tail behind quadratically small logarithmic law error.
The explicit certificates, statistical consequences, and finite checks make
these proposed results suitable for independent review and a staged formal
verification effort. The exact leading constants, the ratio-$1/2$ boundary,
and pairwise stability on separated classes remain distinct next targets.

\begin{thebibliography}{9}
\small\raggedright

\bibitem{predecessor}
Vladimir Reshetnikov and contributors, ProveIt research collection.
\emph{Recovering Uniform Factors: Exact Moment Fibres and Logarithmic
Instability Near the Fabius--Rvachev Law}.
AI-assisted research manuscript, archival arrival of 28 September 2026.
Inspected at repository commit \texttt{\repoSHA}; exact path and blob are
recorded in \cref{sec:provenance}.
\href{https://github.com/VladimirReshetnikov/ProveIt/blob/39472967ce67566d7c530fd5b8d6a3a95e68fe22/Analysis/FabiusFunction/docs/semi-formalized-research-frontiers/drafts/inverse-and-sampling/Recovering_Uniform_Factors_Fabius_Rvachev/article.tex}{Pinned source manuscript}.
The adjacent \href{https://github.com/VladimirReshetnikov/ProveIt/blob/39472967ce67566d7c530fd5b8d6a3a95e68fe22/Analysis/FabiusFunction/docs/semi-formalized-research-frontiers/drafts/inverse-and-sampling/README.md}{collection README}
records its unreviewed status.

\bibitem{billey}
Sara C. Billey and Joshua P. Swanson.
\emph{The metric space of limit laws for $q$-hook formulas}.
Combinatorial Theory \textbf{2}(2), 2022.
DOI: \href{https://doi.org/10.5070/C62257868}{10.5070/C62257868}.
\url{https://arxiv.org/abs/2010.12701v2}.
The generalized-uniform and Gaussian-extended parametrizations are discussed
in Section 3; qualitative identifiability is established background.

\bibitem{arias}
Juan Arias de Reyna.
\emph{An infinitely differentiable function with compact support:
Definition and properties}.
arXiv:1702.05442, 2017.
English translation of ``Definici\'on y estudio de una funci\'on
indefinidamente diferenciable de soporte compacto,''
Rev. Real Acad. Ciencias \textbf{76} (1982), 21--38.
\url{https://arxiv.org/abs/1702.05442}.

\bibitem{massart}
Pascal Massart.
\emph{The tight constant in the Dvoretzky--Kiefer--Wolfowitz inequality}.
Annals of Probability \textbf{18}(3) (1990), 1269--1283.
DOI: \href{https://doi.org/10.1214/aop/1176990746}{10.1214/aop/1176990746}.

\bibitem{weidudley}
Fan Wei and Richard M. Dudley.
\emph{Dvoretzky--Kiefer--Wolfowitz inequalities for the two-sample case}.
arXiv:1107.5356, 2011.
\url{https://arxiv.org/abs/1107.5356}.
The introductory account records the one-sample Massart bound used here;
no two-sample extension is required by our testing argument.

\end{thebibliography}
\end{document}
